% =====================================================================
%  CLASE 10 -- Drenaje transversal de caminos y carreteras: alcantarillas
% =====================================================================
\clase{10}{Drenaje transversal de caminos y carreteras: alcantarillas}

\section{Drenaje transversal de caminos y carreteras}
El objetivo es no interrumpir cursos de agua existentes, por ejemplo: ríos,
quebradas, esteros, etc. Existen distintas soluciones:
\begin{enumerate}
  \item \textbf{Tubos de hormigón o HDPE} (polietileno de alta densidad):
  \begin{itemize}
    \item Si se tiene más de un tubo, la distancia libre entre tubos debe ser $D/2$.
    \item Solución utilizada para caudales pequeños.
    \item Más recomendaciones constructivas en el Manual de Carreteras V4.
  \end{itemize}
  \item \textbf{Cajones de hormigón armado:}
  \begin{itemize}
    \item Recomendaciones constructivas en el Manual de Carreteras V4.
    \item La solución de cajones múltiples compite con puentes pequeños.
  \end{itemize}
  \item \textbf{Puentes} (estribos y cepas): solución utilizada cuando se
        tienen caudales de magnitud importante, por ejemplo ríos o grandes quebradas.
\end{enumerate}
Tanto para cajones como para tubos se deben considerar \textbf{muros guía} a
la entrada y salida. También se conocen como \textbf{atraviesos o alcantarillas}.

\begin{figure}[H]
\centering
\begin{tikzpicture}[font=\small]
  \fill[suelo!80] (0,0) rectangle (9,1.8);
  \fill[gray!25] (0,1.8) .. controls (2.5,1.8) and (3.3,0.4) .. (4.5,0.4) .. controls (5.7,0.4) and (6.5,1.8) .. (9,1.8) -- cycle;
  \draw[thick] (0,1.8) -- (9,1.8);
  \draw[thick,fill=white] (4.5,1.0) circle (0.65);
  \begin{scope} \clip (4.5,1.0) circle (0.65); \fill[aguaclara] (3.8,0.3) rectangle (5.2,1.2); \end{scope}
  \draw[thick] (4.5,1.0) circle (0.65);
  \pic at (4.5,1.22) {nivelagua};
  \draw[dashed] (3.85,1.0) -- (3.85,-0.4) (5.15,1.0) -- (5.15,-0.4);
  \draw[cota] (3.85,-0.3) -- (5.15,-0.3) node[midway,below]{$D$};
  \draw[<-] (6.3,1.4) -- (7.2,2.3) node[above,azulusm]{Terraplén};
  \draw[<-] (5.4,0.7) -- (6.6,-0.3) node[right,azulusm]{Material de relleno};
\end{tikzpicture}
\caption{Alcantarilla de tubo bajo un terraplén (corte transversal al cauce).}
\end{figure}

\section{Alcantarillas}
\begin{itemize}
  \item Son conductos relativamente cortos a través de los cuales se cruza el
        agua por debajo del camino o carretera, desde un lado al otro de éste.
        Se entiende por alcantarilla una \textbf{estructura de drenaje} cuya
        luz, medida paralela al eje de la carretera, sea de \textbf{hasta seis
        metros (ancho)}; losas mayores se tratarán como puentes (MC, 3.703.101).
  \item Se consideran estructuras menores en el diseño y construcción de una
        vía, pero deben ser diseñadas y construidas para soportar las cargas
        del tránsito, el peso de tierra sobre ellas y las cargas durante la construcción.
  \item El diseño consiste en determinar el \textbf{diámetro más económico}
        que permita pasar el caudal de diseño sin exceder la carga máxima de
        entrada y atendiendo criterios de arrastre de sedimentos (velocidades
        de escurrimiento).
  \item Criterios para este tipo de obras:
  \begin{enumerate}[label=\alph*)]
    \item En planta, las obras se dispondrán en la dirección del cauce.
    \item En perfil, se deberá ajustar al perfil de la obra.
    \item La sección debe respetar las dimensiones del cauce natural y no
          provocar fuertes estrechamientos.
  \end{enumerate}
\end{itemize}

\subsection{Algunas recomendaciones (Horacio Mery)}
\begin{itemize}
  \item El nivel aguas arriba sea similar al que se tiene en el cauce para la
        crecida de diseño. ¿Para qué?
  \item El nivel de aguas arriba debe resguardar la seguridad del terraplén.
  \item La mayoría de las alcantarillas se diseñan para operar con
        escurrimiento libre cercano al crítico, logrando $Fr$ entre 0.7 y 0.9.
  \item Se busca maximizar el caudal pasante, reduciendo la sección
        transversal y por tanto el costo.
\end{itemize}

\paragraph{Tragedia del estero Minte (07/05/1995).} El ancho de la
alcantarilla era algo estrecho frente a las precipitaciones (192 mm de agua
en 48 horas); la entrada de la alcantarilla se tapó con árboles y barro,
generando la falla de la autopista, que causó la muerte de 27 personas, 15
de ellas niños. La justicia dictaminó que el Estado fue responsable como
consecuencia de un mal diseño.

\section{Diseño hidráulico de alcantarillas}
\subsection{Recomendaciones de diseño}
\begin{itemize}
  \item La tubería o cajón debe seguir la línea de flujo del agua (en planta,
        alineada con el cauce; no forzar el cauce a cruzar perpendicular al camino).
  \item En un cauce amplio es mejor usar tubos múltiples que dispersen el
        flujo, en lugar de concentrarlo en un solo tubo.
  \item La tubería o cajón debe seguir la pendiente del terreno: no instalarla
        muy profunda ni muy alta; no alterar la elevación del fondo del canal.
\end{itemize}

\begin{figure}[H]
\centering
\begin{tikzpicture}[font=\small]
  \foreach \X/\ok in {0/0,5/1}{
  \begin{scope}[xshift=\X cm]
    \fill[gray!50] (-0.3,1.2) rectangle (3.7,2.0);
    \ifnum\ok=0
      \fill[aguaclara,draw=azulusm] (0.9,3.2) -- (1.5,3.2) -- (1.6,2.0) -- (1.0,2.0) -- cycle;
      \fill[aguaclara,draw=azulusm] (1.0,1.2) -- (1.6,1.2) -- (1.5,0) -- (0.9,0) -- cycle;
      \fill[gray!20,draw] (1.0,1.2) rectangle (1.6,2.0);
      \draw[rojo,line width=3pt] (2.3,0.3) -- (3.1,1.0) (2.3,1.0) -- (3.1,0.3);
    \else
      \fill[aguaclara,draw=azulusm] (1.8,3.2) -- (2.4,3.2) -- (1.9,2.0) -- (1.3,2.0) -- cycle;
      \fill[aguaclara,draw=azulusm] (1.0,1.2) -- (1.6,1.2) -- (1.1,0) -- (0.5,0) -- cycle;
      \fill[gray!20,draw] (1.3,2.0) -- (1.9,2.0) -- (1.6,1.2) -- (1.0,1.2) -- cycle;
      \draw[azulusm,line width=3pt] (2.4,0.7) -- (2.7,0.3) -- (3.4,1.1);
    \fi
  \end{scope}}
\end{tikzpicture}
\caption{La alcantarilla debe seguir la línea de flujo del cauce (derecha), no
forzar un cruce perpendicular (izquierda).}
\end{figure}

% ---------------------------------------------------------------------
% Macro: perfil longitudinal de una alcantarilla
% #1: contenido extra
\newcommand{\perfilalcantarilla}[1]{%
\begin{tikzpicture}[font=\small,scale=1.0]
  % fondo inclinado
  \draw[thick,gray!70!black] (0,0.6) -- (9,0);
  \draw[dashed] (0,0.2) -- (8,0.2);
  % alcantarilla (cajón) y terraplén
  \fill[suelo!80] (3.1,2.6) -- (3.4,2.8) -- (7.6,2.5) -- (7.9,2.0) -- (7.9,1.4) -- (3.1,1.8) -- cycle;
  \draw[thick] (3.0,2.8) -- (3.0,1.85) -- (8.0,1.5) -- (8.0,2.0);
  #1
\end{tikzpicture}}

\subsection{Criterio 1: diseño con control de entrada}
Se diseña la alcantarilla para que el escurrimiento en ella sea
\textbf{torrencial}:
\begin{center}
\colorbox{azulusm}{\parbox{0.7\textwidth}{\centering\color{white}\large
Independencia aguas arriba de lo que ocurre aguas abajo}}
\end{center}
Así, asumiendo la condición más crítica, es decir aguas arriba de la
alcantarilla se tiene escurrimiento de río:
\begin{center}
\colorbox{azulusm}{\parbox{0.7\textwidth}{\centering\color{white}\large
Caída hidráulica en la entrada de la alcantarilla}}
\end{center}

\begin{figure}[H]
\centering
\perfilalcantarilla{
  \fill[aguaclara] (0,0.6) -- (0,2.4) .. controls (2,2.3) and (2.6,1.4) .. (3.1,1.1) .. controls (5,0.7) and (7,0.4) .. (8,0.35) -- (8,0.06) -- cycle;
  \draw[azulusm,thick] (0,2.4) .. controls (2,2.3) and (2.6,1.4) .. (3.1,1.1) .. controls (5,0.7) and (7,0.4) .. (8,0.35);
  \draw[cota] (0.3,0.58) -- (0.3,2.4) node[midway,right]{$H=h_n$};
  \draw[cota] (3.2,0.4) -- (3.2,1.08) node[midway,right]{$h_c$};
  \draw[cota] (7.3,0.12) -- (7.3,1.55) node[midway,left]{$D$};
  \draw[cota] (8.2,0.05) -- (8.2,0.35) node[midway,right]{$h_T$};
  \node[azulusm] at (0.8,2.8) {Río}; \node[azulusm,rotate=-5] at (5.4,0.9) {Torrente};
  \node at (1.5,0.35) {$S_D$};
}
\caption{Control de entrada: escurrimiento de río aguas arriba y torrente en la alcantarilla.}
\end{figure}

\paragraph{Sección rectangular.} Con $H=h+\frac{v^2}{2g}$:
\[
H_c=\frac32h_c,\qquad h_c=\sqrt[3]{\frac{Q^2}{gb^2}}
\]
Despreciando pérdidas de energía, $H_c=H$; y considerando $D=H_c$:
\begin{formula}
$\displaystyle D=\frac32\frac{1}{\sqrt[3]g}\left(\frac Qb\right)^{2/3}$
\end{formula}
Para determinar el ancho $b$ de la alcantarilla se impone $D=H$. Luego,
considerando la ecuación de Manning para $h=h_c=\frac23H$:
\begin{formula}
$\displaystyle S_c=\frac23g\cdot H\cdot n^2\left(\frac{b+\frac43H}{\frac23H\cdot b}\right)^{4/3}$
\end{formula}
La pendiente de diseño debe cumplir con
\[
S_D\geq S_c\ \Rightarrow\ \text{torrente};\qquad\text{si } S_D\leq S_c\ \Rightarrow\ S_D=S_c
\]
Una vez definida la geometría se debe verificar que $v<5$ [m/s].

\paragraph{Sección circular} (considerando $h_c=0.7D$, con $D$ en [m] y $Q$ en [m$^3$/s]):
\begin{formula}
$Q=1.425\,D^{5/2}\qquad\qquad\displaystyle S_c=\frac{31.14\,n^2}{D^{1/3}}$
\end{formula}
Por seguridad se debe diseñar con $D>H$.

\begin{table}[H]
\centering\small
\renewcommand{\arraystretch}{2.0}
\begin{tabular}{l>{$}c<{$}>{$}c<{$}>{$}c<{$}>{$}c<{$}}
\toprule
\textbf{Sección} & \text{Área } A & \text{Perímetro } P & \text{Radio hidráulico } R_h & \text{Ancho superficial } b_s\\
\midrule
Rectangular & by & b+2y & \dfrac{by}{b+2y} & b\\
Trapezoidal & (b+ky)y & b+2y\sqrt{1+k^2} & \dfrac{(b+ky)y}{b+2y\sqrt{1+k^2}} & b+2ky\\
Triangular & ky^2 & 2y\sqrt{1+k^2} & \dfrac{ky}{2\sqrt{1+k^2}} & 2ky\\
Circular & \dfrac{(\theta-\operatorname{sen}\theta)D^2}{8} & \dfrac{\theta D}{2} & \left(1-\dfrac{\operatorname{sen}\theta}{\theta}\right)\dfrac D4 & \operatorname{sen}\left(\dfrac\theta2\right)D\ \text{ó}\ 2\sqrt{y(D-y)}\\
Parabólica & \dfrac23 yb_s & b_s+\dfrac{8y^2}{3b_s} & \dfrac{2b_s^2y}{3b_s^2+8y^2} & \dfrac{3A}{2y}\\
\bottomrule
\end{tabular}
\caption{Elementos geométricos de secciones típicas; en la sección circular
$\theta=2\cos^{-1}\left(1-\frac{2y}{D}\right)$.}
\end{table}

\subsection{Criterio 2: diseño y verificación}
Se toman en cuenta los caudales correspondientes a 2 períodos de retorno:
\begin{itemize}
  \item $T_1$: diseño con control de entrada, p.~ej.\ $T_1=25$ años.
  \item $T_2$: verificación a sección completa, p.~ej.\ $T_2=50$ años
        (entrada ahogada).
\end{itemize}
Se busca que el flujo no cause sobrepresiones en la cara superior del cajón.

\paragraph{Verificación a sección completa.} Se calcula la pendiente normal
para $Q_{2D}$ ($T_2$) imponiendo $h_n=D$:
\begin{formula}
Sección circular: $\displaystyle S_n=\frac{10.3\,Q_{2D}^2\cdot n^2}{D^{16/3}}$\qquad\qquad
Sección rectangular: $\displaystyle\frac{Q_{2D}\cdot n}{\sqrt{S_n}}=\Omega\cdot R_H^{2/3}$
\end{formula}
La comparación con la pendiente de diseño ($S_D$) genera tres casos:
\begin{enumerate}
  \item $S_D>S_n$
  \item $S_D=S_n$
  \item $S_D<S_n$
\end{enumerate}

\paragraph{Caso 1: $S_D>S_n$.} La alcantarilla actúa como compuerta
(contracción a la entrada). Se debe calcular $H'$ para verificar que no se
inunde el camino. Se conserva la energía entre (1) y (2):
\begin{formula}
$\displaystyle H'=c_c\cdot D+\frac{1}{2g}\left(\frac{Q_{2D}}{b\cdot c_c\cdot D}\right)^2,\qquad c_c\approx0.611$
\end{formula}

\begin{figure}[H]
\centering
\perfilalcantarilla{
  \fill[aguaclara] (0,0.6) -- (0,2.7) -- (3.0,2.7) -- (3.0,1.85) .. controls (3.4,1.2) and (4,1.0) .. (5,0.85) .. controls (6.5,0.55) and (7.5,0.4) .. (8,0.35) -- (8,0.06) -- cycle;
  \draw[azulusm,thick] (0,2.7) -- (3.0,2.7);
  \draw[azulusm,thick] (3.0,1.85) .. controls (3.4,1.2) and (4,1.0) .. (5,0.85) .. controls (6.5,0.55) and (7.5,0.4) .. (8,0.35);
  \pic at (2.5,2.72) {nivelagua};
  \draw[cota] (1.3,0.52) -- (1.3,2.7) node[midway,right]{$H'$};
  \draw[cota] (4.3,0.33) -- (4.3,0.95) node[midway,right]{$h_2=c_cD$};
  \draw[cota] (7.3,0.12) -- (7.3,1.55) node[midway,left]{$D$};
  \node at (1.3,0.2) {(1)}; \node at (4.3,0) {(2)};
  \node[azulusm] at (0.6,2.4) {Río}; \node[azulusm,rotate=-5] at (6,0.95) {Torrente};
}
\caption{Caso 1 ($S_D>S_n$): la alcantarilla actúa como compuerta.}
\end{figure}

\paragraph{Caso 2: $S_D=S_n$.} Se puede realizar un balance de energía entre
(1) y (2) considerando la pérdida singular a la entrada:
\begin{formula}
$\displaystyle H'=D+\frac1{2g}\,(k_e+1)\cdot\frac{Q_{2D}^2}{(D\cdot b)^2}$
\end{formula}
El valor de $k_e$ dependerá de la forma del flujo de entrada:
\begin{center}
\begin{tabular}{lc}
\toprule
Tipo de entrada & $k_e$\\
\midrule
Tubo que sobresale del muro (entrada reentrante) & 0.9\\
Entrada abocinada (redondeada) & 0.2\\
Entrada a ras del muro (arista viva) & 0.5\\
\bottomrule
\end{tabular}
\end{center}

\paragraph{Caso 3: $S_D<S_n$.} Se producen sobrepresiones sobre la cara
posterior de la alcantarilla; el escurrimiento aguas arriba se ve
influenciado por lo que ocurre aguas abajo. Se realiza un balance de energía
entre (1) y (2) considerando las pérdidas singulares y por fricción:
\begin{formula}
$\displaystyle H'=D+\frac{v_2^2}{2g}+\Lambda_e+J\cdot L-S_D\cdot L,\qquad
\Lambda_e=k_e\frac{v_2^2}{2g},\qquad J=\frac{v_2^2\cdot n^2}{R_H^{4/3}}$
\end{formula}

\begin{figure}[H]
\centering
\perfilalcantarilla{
  \fill[aguaclara] (0,0.6) -- (0,2.7) -- (3.0,2.7) -- (3.0,1.85) -- (8.0,1.5) -- (8.0,0.0) -- (0,0.6);
  \draw[azulusm,thick] (0,2.7) -- (3.0,2.7);
  \pic at (2.5,2.72) {nivelagua};
  \foreach \y in {0.3,0.6,0.9,1.2} \draw[->,azulusm] (8.0,\y) .. controls (8.4,\y) .. (8.8,\y+0.2);
  \draw[cota] (1.3,0.52) -- (1.3,2.7) node[midway,right]{$H'$};
  \draw[cota] (7.3,0.06) -- (7.3,1.55) node[midway,left]{$D$};
  \draw[dashed] (3.0,0.37) -- (3.0,-0.6) (8.0,0.0) -- (8.0,-0.6);
  \draw[cota] (3.0,-0.5) -- (8.0,-0.5) node[midway,below]{$L$};
  \node at (1.3,0.2) {(1)}; \node[right] at (8.0,-0.2) {(2)};
}
\caption{Caso 3 ($S_D<S_n$): alcantarilla trabajando a sección llena.}
\end{figure}

\paragraph{Adicional: control de salida.} La pendiente aguas abajo es tan
pequeña que el escurrimiento de río ahoga al atravieso. Se realiza un
balance de energía entre (1) y (3) considerando las pérdidas singulares
(contracción a la entrada y expansión a la salida) y por fricción:
\begin{formula}
$\displaystyle H'=D+\frac{v_2^2}{2g}+\Lambda_e+\Lambda_s+J\cdot L-S_D\cdot L$
\end{formula}
donde
\[
\Lambda_s=k_s\frac{(v_2-v_3)^2}{2g},\qquad\Lambda_e=k_e\frac{v_2^2}{2g},\qquad J=\frac{v_2^2\cdot n^2}{R_H^{4/3}},
\]
\[
k_s=0.22-0.44\ \ (k_s=0.4\ \text{recomendado})
\]
y $h_3=h_{n3}$ es la altura normal aguas abajo. Referencia: FHWA, \emph{Hydraulic
Design of Highway Culverts},
\url{https://www.fhwa.dot.gov/engineering/hydraulics/library_arc.cfm?id=13&pub_number=7}.

\subsection{Períodos de retorno}
\begin{table}[H]
\centering\small
\resizebox{\linewidth}{!}{%
\begin{tabular}{llcccccc}
\toprule
& & \multicolumn{2}{c}{$T$ [años]} & Vida útil & \multicolumn{2}{c}{Riesgo de falla [\%]}\\
\cmidrule(lr){3-4}\cmidrule(lr){6-7}
\textbf{Tipo de obra} & \textbf{Ruta} & Diseño & Verif. & $n$ [años] & Diseño & Verif.\\
\midrule
Puentes y viaductos & Carreteras & 200 & 300 & 50 & 22 & 15\\
 & Caminos & 100 & 150 & 50 & 40 & 28\\
Alcantarillas ($S>1.75$ m$^2$) o $H$ terrap.\ $\geq10$ m & Carreteras & 100 & 150 & 50 & 40 & 28\\
\quad y estructuras enterradas & Caminos & 50 & 100 & 30 & 45 & 26\\
Alcantarillas ($S<1.75$ m$^2$) & Carreteras & 50 & 100 & 50 & 64 & 40\\
 & Caminos & 25 & 50 & 30 & 71 & 45\\
\bottomrule
\end{tabular}}
\caption{Períodos de retorno (Manual de Carreteras, Volumen 3). $S$ es el área útil de la alcantarilla.}
\end{table}

\subsection{Alturas y velocidades}
\begin{table}[H]
\centering\small
\begin{tabular}{llll}
\toprule
\textbf{Tipo de cauce} & \textbf{Tubos} & \textbf{Cajones} & \textbf{Losas} ($L\leq6.0$ m)$^*$\\
\midrule
Canales & $D$ (diámetro) & $H$ (altura total) & $H-0.10$ m\\
Diseño cauces naturales & $D+0.3$ m & $H+0.3$ m & $H-0.10$ m\\
Verificación cauces naturales & $D+0.6$ m & $H+0.6$ m & $H$\\
\multicolumn{4}{l}{\quad Pero $H_e$ máximo no puede sobrepasar la cota exterior del SAP $-0.3$ m}\\
\bottomrule
\end{tabular}
\caption{Tabla 3.703.301.A: carga hidráulica de diseño ($H_e$, m).
$^*$Si $L>6.0$ m, revancha como en puentes.}
\end{table}

\begin{table}[H]
\centering\small
\begin{tabular}{lcc}
\toprule
\textbf{Tipo de terreno} & \textbf{Flujo intermitente} [m/s] & \textbf{Flujo permanente} [m/s]\\
\midrule
Arena fina (no coloidal) & 0.75 & 0.75\\
Arcilla arenosa (no coloidal) & 0.75 & 0.75\\
Arcilla limosa (no coloidal) & 0.90 & 0.90\\
Arcilla fina & 1.00 & 1.00\\
Ceniza volcánica & 1.20 & 1.00\\
Grava fina & 1.50 & 1.20\\
Arcilla dura (coloidal) & 1.80 & 1.40\\
\multicolumn{3}{l}{\emph{Material graduado (no coloidal):}}\\
\quad desde arcilla a grava & 2.00 & 1.50\\
\quad desde limo a grava & 2.10 & 1.70\\
Grava & 2.30 & 1.80\\
Grava gruesa & 2.40 & 2.00\\
Desde grava a piedras (bajo 15 cm) & 2.70 & 2.10\\
Desde grava a piedras (sobre 20 cm) & 3.00 & 2.40\\
\bottomrule
\end{tabular}
\caption{Tabla 3.703.301.B: velocidades máximas admisibles en canales no
revestidos (fuente: Manual de Carreteras de California).}
\end{table}
